% GENERATED FILE. Edit canonical lessons, then run npm run generate:books.
% canonical-source-hash: fnv1a64:d54d03f269450e1b
% canonical-lessons: ZH-C217-meiguanxi, ZH-C217-dangran, ZH-C217-jide, ZH-C217-dating, ZH-C217-jiejue, ZH-R217-review-212-217, ZH-R217-review-final-217

\chapter{Never Mind}
\label{ch:zh-never-mind}
\hlchaptermodality{217}

\begin{chapteropening}
\textbf{By the end of this chapter:} \emph{I can say it doesn't matter, say of course, say I remember, ask around and solve a problem.}

Complete the last lesson of chapter 217: I can say it doesn't matter, say of course, say I remember, ask around and solve a problem.
\end{chapteropening}

\section[méi guānxi]{méi guānxi — it doesn't matter; never mind}

\label{lesson:ZH-C217-meiguanxi}



\begin{quote}
Before the new one: say the Mandarin for police officer, then the Mandarin for to call the police.
\end{quote}

\subsection*{You'll want to know: méi guānxi}
\emph{méi guānxi} — \textquotedblleft{}it doesn't matter; never mind\textquotedblright{}.

\begin{quote}
\emph{méi guānxi wǒmen míngtiān qù} — Never mind, we'll go tomorrow.
\end{quote}

Its characters come later; for now it is a word for the ear.

\begin{cousinweb}
\emph{méi guānxi} is built from \zh{没} 'not have' and \emph{guānxi} 'a connection', all three parts inherited from Middle Chinese, \emph{mwot}, \emph{kwaen} and \emph{kejH} (Guangyun, in Baxter's transcription): no connection to worry about.
\end{cousinweb}

\begin{grammarlens}[title={What you've built}]
One more word for saying no and asking.
\end{grammarlens}

\subsection*{Your turn}
\begin{itemize}
\raggedright
  \item \emph{Say it:} \emph{méi guānxi}
  \item \emph{Say it:} \emph{méi guānxi}, once more
  \item \emph{Say it:} \emph{méi guānxi wǒmen míngtiān qù} — Never mind, we'll go tomorrow.
  \item \emph{Recall:} say the Mandarin for police officer, then the Mandarin for to call the police, then say \emph{méi guānxi} again
\end{itemize}

\begin{tcolorbox}[breakable,colback=teal!4,colframe=teal!35!black,title={Before you move on}]
What does méi guānxi mean? (\textquotedblleft{}It doesn't matter; never mind\textquotedblright{}.) Say it once more.
\end{tcolorbox}
\section[dāngrán]{dāngrán — of course}

\label{lesson:ZH-C217-dangran}



\begin{quote}
Before the new one: say the Mandarin for to call the police, then the Mandarin for never mind.
\end{quote}

\subsection*{You'll want to know: dāngrán}
\emph{dāngrán} — \textquotedblleft{}of course\textquotedblright{}.

\begin{quote}
\emph{dāngrán kěyǐ} — Of course you can.
\end{quote}

Its characters come later; for now it is a word for the ear.

\begin{cousinweb}
\emph{dāngrán} is built from \emph{dāng} 'to match' and \emph{rán} 'so', both inherited from Middle Chinese, \emph{tang} and \emph{nyen} (Guangyun, in Baxter's transcription), which glosses the second 'like this' and 'it is so'.
\end{cousinweb}

\begin{grammarlens}[title={What you've built}]
One more word for saying no and asking.
\end{grammarlens}

\subsection*{Your turn}
\begin{itemize}
\raggedright
  \item \emph{Say it:} \emph{dāngrán}
  \item \emph{Say it:} \emph{dāngrán}, once more
  \item \emph{Say it:} \emph{dāngrán kěyǐ} — Of course you can.
  \item \emph{Recall:} say the Mandarin for to call the police, then the Mandarin for never mind, then say \emph{dāngrán} again
\end{itemize}

\begin{tcolorbox}[breakable,colback=teal!4,colframe=teal!35!black,title={Before you move on}]
What does dāngrán mean? (\textquotedblleft{}Of course\textquotedblright{}.) Say it once more.
\end{tcolorbox}
\section[jìde]{jìde — to remember}

\label{lesson:ZH-C217-jide}



\begin{quote}
Before the new one: say the Mandarin for never mind, then the Mandarin for of course.
\end{quote}

\subsection*{You'll want to know: jìde}
\emph{jìde} — \textquotedblleft{}to remember\textquotedblright{}.

\begin{quote}
\emph{nǐ hái jìde wǒ ma} — Do you still remember me?
\end{quote}

Its characters come later; for now it is a word for the ear.

\begin{cousinweb}
\emph{jìde} is built from \emph{jì} 'to note down' and \zh{得} 'to get', both inherited from Middle Chinese, \emph{kiH} and \emph{tok} (Guangyun, in Baxter's transcription); in the word the second loses its tone.
\end{cousinweb}

\begin{grammarlens}[title={What you've built}]
One more word for saying no and asking.
\end{grammarlens}

\subsection*{Your turn}
\begin{itemize}
\raggedright
  \item \emph{Say it:} \emph{jìde}
  \item \emph{Say it:} \emph{jìde}, once more
  \item \emph{Say it:} \emph{nǐ hái jìde wǒ ma} — Do you still remember me?
  \item \emph{Recall:} say the Mandarin for never mind, then the Mandarin for of course, then say \emph{jìde} again
\end{itemize}

\begin{tcolorbox}[breakable,colback=teal!4,colframe=teal!35!black,title={Before you move on}]
What does jìde mean? (\textquotedblleft{}To remember\textquotedblright{}.) Say it once more.
\end{tcolorbox}
\section[dǎting]{dǎting — to ask around; to inquire}

\label{lesson:ZH-C217-dating}



\begin{quote}
Before the new one: say the Mandarin for of course, then the Mandarin for to remember.
\end{quote}

\subsection*{You'll want to know: dǎting}
\emph{dǎting} — \textquotedblleft{}to ask around; to inquire\textquotedblright{}.

\begin{quote}
\emph{wǒ dǎtingle shūdiàn jīntiān bù kāimén} — I asked around, and the bookshop is closed today.
\end{quote}

Its characters come later; for now it is a word for the ear.

\begin{cousinweb}
\emph{dǎting} is built from \zh{打} 'to strike' and \emph{tīng} 'to listen', both inherited from Middle Chinese, \emph{taengX} and \emph{theng} (Guangyun, in Baxter's transcription); in the word the second loses its tone.
\end{cousinweb}

\begin{grammarlens}[title={What you've built}]
One more word for saying no and asking.
\end{grammarlens}

\subsection*{Your turn}
\begin{itemize}
\raggedright
  \item \emph{Say it:} \emph{dǎting}
  \item \emph{Say it:} \emph{dǎting}, once more
  \item \emph{Say it:} \emph{wǒ dǎtingle shūdiàn jīntiān bù kāimén} — I asked around, and the bookshop is closed today.
  \item \emph{Recall:} say the Mandarin for of course, then the Mandarin for to remember, then say \emph{dǎting} again
\end{itemize}

\begin{tcolorbox}[breakable,colback=teal!4,colframe=teal!35!black,title={Before you move on}]
What does dǎting mean? (\textquotedblleft{}To ask around; to inquire\textquotedblright{}.) Say it once more.
\end{tcolorbox}
\section[jiějué]{jiějué — to solve; to settle}

\label{lesson:ZH-C217-jiejue}



\begin{quote}
Before the new one: say the Mandarin for to remember, then the Mandarin for to ask around.
\end{quote}

\subsection*{You'll want to know: jiějué}
\emph{jiějué} — \textquotedblleft{}to solve; to settle\textquotedblright{}.

\begin{quote}
\emph{zhège wèntí wǒmen néng jiějué} — We can sort this problem out.
\end{quote}

Its characters come later; for now it is a word for the ear.

\begin{cousinweb}
\emph{jiějué} is built from \emph{jiě} 'to untie' and \emph{jué} 'to break through', both inherited from Middle Chinese, \emph{keaX} and \emph{kwet} (Guangyun, in Baxter's transcription), which glosses the first 'to explain' and 'to loosen', and the second 'to break'.
\end{cousinweb}

\begin{grammarlens}[title={What you've built}]
One more word for saying no and asking.
\end{grammarlens}

\subsection*{Your turn}
\begin{itemize}
\raggedright
  \item \emph{Say it:} \emph{jiějué}
  \item \emph{Say it:} \emph{jiějué}, once more
  \item \emph{Say it:} \emph{zhège wèntí wǒmen néng jiějué} — We can sort this problem out.
  \item \emph{Recall:} say the Mandarin for to remember, then the Mandarin for to ask around, then say \emph{jiějué} again
\end{itemize}

\begin{tcolorbox}[breakable,colback=teal!4,colframe=teal!35!black,title={Before you move on}]
What does jiějué mean? (\textquotedblleft{}To solve; to settle\textquotedblright{}.) Say it once more.
\end{tcolorbox}
\section[(dialogue)]{Looking back: from Services in Town to Never Mind}

\label{lesson:ZH-R217-review-212-217}



\begin{quote}
Say the Mandarin for to print out and to photocopy.
\end{quote}

\subsection*{Your turn}
Say these aloud:

\begin{itemize}
\raggedright
  \item to print out, to photocopy, to top up, to make an appointment, to get a haircut
  \item visa, tour guide, tourist, admission ticket, have a safe journey
  \item problem, to tell someone, to grasp a meaning, a way to deal with it, what should I do
  \item to deliver goods, to sign for, to leave in safekeeping, laundry shop, to dry-clean
  \item oh no, odd, to steal, police officer, to call the police
  \item never mind, of course, to remember, to ask around, to solve
\end{itemize}

\begin{tcolorbox}[breakable,colback=teal!4,colframe=teal!35!black,title={Before you move on}]
To print out? (\emph{dǎyìn}.) To photocopy? (\emph{fùyìn}.) To top up? (\emph{chōngzhí}.) To make an appointment? (\emph{yùyuē}.) To get a haircut? (\emph{lǐfà}.) Visa? (\emph{qiānzhèng}.) Tour guide? (\emph{dǎoyóu}.) Tourist? (\emph{yóukè}.) Admission ticket? (\textbf{\zh{门票}}, \emph{ménpiào}.) Have a safe journey? (\emph{yílù píng'ān}.) Problem? (\emph{wèntí}.) To tell someone? (\emph{gàosu}.) To grasp a meaning? (\textbf{\zh{明白}}, \emph{míngbai}.) A way to deal with it? (\emph{bànfǎ}.) What should I do? (\emph{zěnme bàn}.) To deliver goods? (\emph{sònghuò}.) To sign for? (\emph{qiānshōu}.) To leave in safekeeping? (\emph{jìcún}.) Laundry shop? (\textbf{\zh{洗衣店}}, \emph{xǐyīdiàn}.) To dry-clean? (\emph{gānxǐ}.) Oh no? (\emph{zāogāo}.) Odd? (\emph{qíguài}.) To steal? (\emph{tōu}.) Police officer? (\emph{jǐngchá}.) To call the police? (\emph{bàojǐng}.) Never mind? (\emph{méi guānxi}.) Of course? (\emph{dāngrán}.) To remember? (\emph{jìde}.) To ask around? (\emph{dǎting}.) To solve? (\emph{jiějué}.)
\end{tcolorbox}
\section[(dialogue)]{Never Mind — once more}

\label{lesson:ZH-R217-review-final-217}



\begin{quote}
Say the Mandarin for never mind and of course.
\end{quote}

\subsection*{Your turn}
\begin{itemize}
\raggedright
  \item \emph{Say it:} never mind, of course, to remember, to ask around, to solve
\end{itemize}

\begin{tcolorbox}[breakable,colback=teal!4,colframe=teal!35!black,title={Before you move on}]
Never mind? (\emph{méi guānxi}.) Of course? (\emph{dāngrán}.) To remember? (\emph{jìde}.) To ask around? (\emph{dǎting}.) To solve? (\emph{jiějué}.)
\end{tcolorbox}
